\section{Problem Statement and Design Goals}

The platform addresses the following problem:
\begin{quote}
\emph{How can a Kubernetes-based control plane reliably manage secure enrollment and WebAssembly application life-cycle operations for heterogeneous embedded devices while preserving observability, reproducibility, and declarative status consistency?}
\end{quote}

This problem statement contains three distinct tensions. The first is a resource-model tension: Kubernetes assumes a declarative control plane with explicit status ownership, whereas embedded deployments often rely on imperative, device-local procedures that are opaque to the cloud. The second is a transport tension: constrained devices need lightweight and secure communication paths, but the cloud side requires rich enough semantics to express enrollment, liveness, and application control. The third is an operational tension: even when protocol messages are correct, the platform can still fail if reconciliation logic misinterprets transient conditions and reports an incorrect device state.

To make these tensions actionable, we derive the following design goals.

\begin{itemize}
    \item \textbf{G1 - Unified control plane.} Desired and observed states for devices, applications, and gateways must remain represented in Kubernetes resources so that operators can reason about the platform through a single authoritative interface.
    \item \textbf{G2 - Secure edge attachment.} Device enrollment must occur over a protected channel and must bind transport-level identity to platform-level resource identity in a way that is auditable and robust.
    \item \textbf{G3 - Runtime portability.} The application format and life-cycle interface should remain portable across cloud, fog, and constrained edge targets, motivating the adoption of WebAssembly as a common execution abstraction.
    \item \textbf{G4 - Operational observability.} The platform must expose enough runtime evidence to confirm enrollment, liveness, deployment status, and gateway association without requiring manual inspection of device internals.
    \item \textbf{G5 - Reproducible validation.} The complete stack should be deployable and testable in a commodity research environment so that experiments can be repeated, reviewed, and extended.
\end{itemize}

These goals imply an architectural consequence: the system cannot treat communication, orchestration, and status management as separate concerns. A platform that satisfies G2 but violates G1 or G4 would remain difficult to operate, while a platform that satisfies G1 but not G2 would be insecure by construction. RETROSPECT therefore adopts a deliberately integrated design in which resource modeling, gateway behavior, and controller reconciliation are co-designed rather than layered opportunistically.
